\documentclass[aspectratio=169]{beamer}

\usetheme{Madrid}

\usepackage{amsmath}
\usepackage{amssymb}
\usepackage{graphicx}
\usepackage{booktabs}
\usepackage{listings}

\usepackage{fontspec}
\usepackage{xeCJK}

\setmainfont{Times New Roman}
\setCJKmainfont{Hiragino Sans}

\title{なぜベンチャーキャピタルは存在するのか}
\subtitle{Theory and Canadian Evidence を中心とした分析}
\author{口木}
\date{\today}

\AtBeginSection[]{
\begin{frame}{現在位置}
\tableofcontents[currentsection]
\end{frame}
}

\lstset{
language=Python,
basicstyle=\ttfamily\tiny,
breaklines=true,
frame=single
}

\begin{document}

\begin{frame}\titlepage\end{frame}

\begin{frame}{目次}
\tableofcontents
\end{frame}

\section{問題設定}

\begin{frame}{研究課題}
研究課題：
\begin{block}{Research Question}
なぜ銀行ではなくVCが存在するのか？
\end{block}

本発表：
\begin{enumerate}
\item 情報の非対称性
\item 契約理論
\item モニタリング
\item モンテカルロシミュレーション
\item 実証研究
\end{enumerate}
\end{frame}

\begin{frame}{VCとは何か}
VCの特徴：
\begin{itemize}
\item 高リスク企業へ投資
\item 株式による資金供給
\item 経営支援
\item IPO/M\&Aによる回収
\end{itemize}
\end{frame}

\begin{frame}{銀行との違い}
銀行:
\[
返済能力 \rightarrow 融資
\]

VC: 
\[
将来価値 \rightarrow 出資
\]

\end{frame}

\section{理論背景}

\begin{frame}{情報の非対称性}
情報量：

\[
I_E>I_I
\]

\begin{itemize}
\item 起業家は企業内部情報を知る
\item 投資家は完全には観察できない
\end{itemize}
\end{frame}

\begin{frame}{逆選択モデル}
企業タイプ

\[
\theta\in\{\theta_H,\theta_L\}
\]

確率

\[
P(\theta_H)=p
\]

期待価値：

\[
E(V)=pV_H+(1-p)V_L
\]

\end{frame}

\begin{frame}{逆選択の問題}
平均価値：

\[
\bar V
=
pV_H+(1-p)V_L
\]

\begin{alertblock}{問題}
高品質企業が過小評価される
\end{alertblock}
\end{frame}

\begin{frame}{モラルハザード}
努力：

\[
e
\]

効用：

\[
U=V(e)-C(e)
\]

一階条件：

\[
V'(e)-C'(e)=0
\]

\end{frame}

\begin{frame}{限界条件の解釈}

\[
V'(e)=C'(e)
\]

左辺：努力追加による利益

右辺：努力追加によるコスト

\vspace{0.3cm}

均衡：

利益増加＝コスト増加

\end{frame}

\begin{frame}{VCモニタリング}

監視水準：

\[
m
\]

成功確率：

\[
p(m)
\]

期待収益：

\[
E(R)
=
p(m)R-M(m)
\]

\end{frame}

\begin{frame}{最適モニタリング}

最大化問題：

\[
\max_m p(m)R-M(m)
\]

FOC：

\[
p'(m)R=M'(m)
\]

\end{frame}

\begin{frame}{ステージファイナンス}

第1期：

\[
I_1
\]

第2期：

\[
I_2
\]

継続条件：

\[
V(I_2)>I_2
\]

\end{frame}

\begin{frame}{ベイズ更新}

事前確率：

\[
P(\theta_H)
\]

事後確率：

\[
P(\theta_H|s)
\]

ベイズ則：

\[
P(\theta_H|s)
=
\frac{P(s|\theta_H)P(\theta_H)}
{P(s)}
\]

\end{frame}

\section{論文の実証}

\begin{frame}{論文の主張}

VCの役割：

\begin{itemize}
\item 資金供給だけではない
\item 情報問題の解決
\item 企業成長支援
\end{itemize}

\end{frame}

\begin{frame}{Canadian Evidence}

図挿入:

\begin{center}
\includegraphics[width=.7\textwidth]{canadian_vc.png}
\end{center}

\end{frame}

\begin{frame}{論文の結論}

\begin{enumerate}
\item VCは情報仲介者
\item モニタリングが重要
\item 段階投資が有効
\end{enumerate}

\end{frame}

\section{モンテカルロ}

\begin{frame}{収益モデル}

\[
R=
\begin{cases}
0&(0.90)\\
2&(0.09)\\
100&(0.01)
\end{cases}
\]

\end{frame}








\begin{frame}{結果図}

結果ヒストグラム
（後で挿入）

%\includegraphics[width=.8\textwidth]{histogram.png}

\end{frame}

\begin{frame}{Power Law}

\[
P(X>x)
=
Cx^{-\alpha}
\]

GoogleやUberのような
少数の大成功が全体利益を支える

\end{frame}

\section{結論}

\begin{frame}{結論}

\begin{itemize}
\item VCは情報問題を解決する制度
\item モニタリングが重要
\item 分散投資が必要
\item 少数成功が利益を支える
\end{itemize}

\end{frame}

\begin{frame}{今後の課題}

\begin{itemize}
\item 実データ推定
\item IPOデータ利用
\item 国別比較
\end{itemize}

\end{frame}

\end{document}

\end{frame}

\end{document}